\section{本章小结}
\label{sec:pdeprob-summary}

一个偏微分方程问题由四样东西共同构成：未知场和它所在的区域，方程本身，初值和边界数据，以及“解”的含义和相应的函数空间。只写一条方程不等于提出了问题，$\partial_tu=\partial_x^{2}u$有无穷多解，它们的差别要由数据来确定。

未知量从向量变成函数后，方程可以由局部守恒和具体的通量规律得出，Poisson、输运、热和波动四个模型都写成同一组变量。空间离散后，热方程就是常微分方程组，但矩阵的谱范数按$4/h^{2}$增长，所以不能把常微分方程的存在性定理直接搬过来。

需要什么数据由方程决定。时间阶数决定初始函数的个数，边界条件可以规定场值、通量或两者的组合，输运方程只在入流边界给数据。初值与边值在角点要相容，相容的阶数随所要求的光滑性增加；Neumann问题还有积分约束。不是数据越多越好。

对二阶线性方程，系数矩阵特征值的符号结构决定类型：椭圆型受整个边界约束，抛物型立刻抹平，双曲型以有限速度传播不光滑性。类型在坐标变换下不变，但是逐点定义，也不记录时间的方向。一阶方程不在这个分类中，它们按特征线传播，非线性时特征线会相交。

数据粗糙或演化产生间断时，要改变方程成立的方式。弱解把导数转移到测试函数上；守恒律的弱解可以不唯一，要由熵解选择；Hamilton--Jacobi方程几乎处处意义下有太多解，要由黏性解选择。不同方程需要不同的解的概念，光滑时它们都回到经典解。

存在、唯一和连续依赖是三个要分别回答的问题，答案依赖所选的空间。热方程正向适定、反向不适定，Laplace方程配Dirichlet数据适定、配Cauchy数据不适定，公式形式上都写得出来，所以写出解不等于证明问题适定。本章的结论都是有限时间内的。

后面几章在具体模型中落实这些条件：第~\ref{chap:elliptic-equations}章建立椭圆方程的弱形式和适定性，第~\ref{chap:diffusion-equations}章和第~\ref{chap:wave-equations}章分别讨论抛物型和双曲型的演化，第~\ref{chap:transport-equations}章、第~\ref{chap:scalar-conservation-laws}章和第~\ref{chap:hamilton-jacobi-equations}章处理一阶方程及其解的选择，数值稳定性留到第~\ref{chap:pde-numerical-methods}章的离散问题中讨论。
